\title{\textbf{Monografía}}
\author{\small Máximo D. Utrera, Facultad de Ingeniería, UBA, Buenos Aires, Argentina \hfill \small \today}
\date{} % Suppress autogenerated date

\twocolumn[
\begin{@twocolumnfalse}
  \maketitle
  \thispagestyle{fancy}
\end{@twocolumnfalse}
]

\section*{Resumen}

En este trabajo se exploran diversas aplicaciones de las redes auto-organizadas de Kohonen (SOM). En primer lugar, se diseña una red de Kohonen bidimensional que aprende una distribución uniforme dentro del círculo unitario, verificando la preservación de la topología y extendiendo el análisis a otras figuras geométricas. Luego, se aplica una red SOM al problema del viajero (Traveling Salesman Problem) con 200 o mas ciudades, obteniendo una solución aproximada eficiente mediante la organización de los nodos en un anillo. Finalmente, se emplea una red de Kohonen sobre un conjunto de datos de alta dimensionalidad para su reducción dimensional y análisis clusters. Se visualiza la matriz U para detectar estos últimos, complementando el estudio con métodos adicionales (a modo de bonus) como el análisis de tendencias y la aplicación de K-means para estimar el número óptimo de agrupamientos. Los resultados confirman la capacidad de las SOM para representar topologías complejas y revelar estructuras latentes en los datos.
\vspace{0.25cm}

\section{Preservación de la topología en distribuciones uniformes sobre formas geométricas}

En esta sección se implementa una red auto-organizada de kohonen de dos dimensiones, en forma de grilla en \textit{python} y luego se la entrena con una series de sets de datos de también dos dimensiones, para evaluar la capacidad de esta red de preservar la topología.

\subsection{Implementación}

El algoritmo empleado es el descrito a continuación:

\begin{itemize}
    \item Se crea una matriz de pesos de tamaño $grid\space size \times grid\space size$  y se los inicializa uniformemente entre 0 y 1.
    \item Se itera por una cantidad de épocas y dentro de las mismas se itera por el dataset tomando muestras de manera aleatoria para cada iteración.
    \item Se calculan las distancias euclidianas de cada neurona hacia la entrada aleatoriamente seleccionada y se selecciona como ganadora a la neurona mas cercana.
    \item Se calculan las distancias de las demás neuronas hacia la ganadora y se aplica la función de vecindad, la cual permite la activación (mas tenue) de las neuronas vecinas y la no activación de neuronas lejanas.
    \item Se actualizan los pesos sinápticos y se disminuye el 'radio' de excitación de neuronas vecinas ($\sigma$) con un factor de decaimiento $\alpha_{\sigma}$.
\end{itemize}

Fórmulas utilizadas:

\begin{itemize}
    \item $\text{winner} = \arg\min_{i} \left\| \mathbf{w}_i - \mathbf{x} \right\|_2$
    \item $d_{grid}(i, j) = \left\| \mathbf{r}_i - \mathbf{r}_j \right\|_2$
    \item $\Lambda_i = \exp\left( -\frac{ \left\| \mathbf{r}_i - \mathbf{r}_{\text{winner}} \right\|^2 }{2\sigma^2} \right)$
    \item $\Delta\mathbf{w}_i = \eta \cdot \Lambda_i \cdot (\mathbf{x} - \mathbf{w}_i)$
    \item $\sigma = \sigma \cdot \alpha_{\sigma}$
\end{itemize}

Representación gráfica de la red de kohonen implementada, donde se tiene un plano bidimensional de neuronas y las entradas están directamente conectadas con todas las neuronas (en el caso de este trabajo $sizeY$ y $sizeX$ serán iguales).

\begin{figure}[H]
\centering
\includegraphics[width=0.49\textwidth]{img/kohonen_net.png}
\caption{Red de kohonen 2D}
\end{figure}

\subsection{Figuras a aprender}

Como se menciono anteriormente en esta sección se entrenaran tres redes, una para cada conjunto de datos conteniendo puntos de las siguientes figuras geométricas.

\begin{figure}[H]
\centering
\includegraphics[width=0.49\textwidth]{img/shapes.png}
\caption{Figuras para el entrenamiento}
\end{figure}

Para las primeras dos figuras se generaron $1000$ puntos y para la tercera (pentágono) se generaron $2000$ de los mismos. Siempre utilizando un muestreo con distribución uniforme.

\subsection{Circulo unitario}

La primera red se entreno con los datos del circulo unitario, y contó con las siguientes características:

\begin{itemize}
    \item $\#Neuronas = 7$
    \item $Epochs = 4$
    \item $Cte. Aprendizaje = 0.1$
    \item $\sigma_{inicial} = 3.0; \alpha=0.999$
\end{itemize}

Como se observa a continuación, la red logró preservar correctamente la forma del círculo gracias al proceso de auto-organización. Las neuronas se distribuyeron de manera que cada una quedó cercana a sus vecinas, lo que indica que se activaron de forma conjunta durante el entrenamiento.

\begin{figure}[H]
\centering
\includegraphics[width=0.42\textwidth]{img/circle.png}
\caption{Preservación de la topología circular}
\end{figure}

\subsection{Cuadrado unitario y pentágono}

Para confirmar la capacidad de este tipo de red de preservar otro tipo de formas de datos, se le muestran los dos restantes conjuntos de datos mencionados (el cuadrado y el pentágono) y se obtuvieron los resultados graficados a continuación.

Para el cuadrado se utilizaron los mismos parámetros de la red anterior, sin embargo, eso no dio buenos resultados para el caso del pentágono, es por esto que para este se incrementaron las épocas ($50$) y el radio de vecindad ($7.0$).

\begin{figure}[H]
\centering
\includegraphics[width=0.42\textwidth]{img/square.png}
\caption{Preservación de la topología cuadrada}
\end{figure}

\begin{figure}[H]
\centering
\includegraphics[width=0.42\textwidth]{img/pentagon.png}
\caption{Preservación de la topología pentagonal}
\end{figure}

Ambas redes exitosamente preservaron la topología de las figuras por medio del ordenamiento de neuronas y vecindades.

\newpage

\section{Problema del viajante de comercio}

En esta sección se aborda el uso de redes de kohonen para la resolución (aproximada) de un problema NP-hard (TSP) utilizando una técnica demostrada en \textit{Introduction To The Theory Of Neural Computation}.

\subsection{Explicación del problema}

El problema del viajante de comercio consiste en encontrar el recorrido más corto posible que permite a un viajante visitar una serie de ciudades una única vez y regresar a la ciudad de origen. Este problema pertenece a la clase de problemas NP-hard, lo que implica que no existen algoritmos eficientes conocidos que lo resuelvan de forma exacta en un tiempo razonable para grandes cantidades de ciudades.

\subsection{Implementación}

En el contexto de redes neuronales, y en particular mediante el uso de redes de Kohonen, el TSP puede abordarse como una tarea de mapeo auto-organizado. Se emplea una técnica que consiste en la creación de un anillo de neuronas que, a través del proceso de entrenamiento no supervisado, se adapta para formar una trayectoria cerrada que recorra todas las ciudades minimizando la distancia total. Durante el entrenamiento, las neuronas se 'mueven' hacia las ciudades y se reorganizan de manera de preservar una estructura continua, lo que permite aproximar una solución al problema.

Es por esto que para esta sección se implemento una red de kohonen uni-dimencional y circular, cuyos únicos cambios se dan en la asignación de pesos y la función de vecindad que ahora considera distancias circulares.

\subsection{Resolución del problema para 200 ciudades}

Se comienza por la generación de puntos aleatorios y uniformes $x,y \in [0, 1]$ para el conjunto de datos de ciudades. Luego se inicializa la red, con $2$ entradas y $3000$ neuronas, teniendo en cuenta lo mencionado anteriormente sobre el anillo de neuronas (en este caso de radio $0.5$), de manera que se obtiene el siguiente gráfico inicial.

\begin{figure}[H]
\centering
\includegraphics[width=0.42\textwidth]{img/tsp_init.png}
\caption{Estado inicial de red de kohonen en anillo}
\end{figure}

Finalmente se lleva a cabo el entrenamiento durante $20$ épocas, un $\sigma_{inicial}=15.0$ y una constante de aprendizaje de $0.8$, y obteniendo el resultado graficado a continuación en un total de $2.1seg$.

\begin{figure}[H]
\centering
\includegraphics[width=0.42\textwidth]{img/tsp_sol.png}
\caption{Solución de TSP con 200 ciudades}
\end{figure}

La red logro efectivamente pasar de un estado inicial de anillo a una solución aproximada del problema que pasa por todas las ciudades del conjunto de datos minimizando el camino entre cada una de ellas y conformando un camino cerrado.

\subsection{Resolución del problema para mas ciudades}

Posterior a la resolución de TSP con 200 ciudades se realizaron dos intentos mas, poniendo a prueba la red con una cantidad medianamente mayor de ciudades y luego con un incremento de $5$ veces la cantidad de ciudades manteniendo las mismas características anteriormente mencionadas de la red. Ambos resultados se grafican a continuación.

\begin{figure}[H]
\centering
\includegraphics[width=0.42\textwidth]{img/tsp_sol_300.png}
\caption{Solución de TSP con 300 ciudades}
\end{figure}

\begin{figure}[H]
\centering
\includegraphics[width=0.42\textwidth]{img/tsp_sol_1000.png}
\caption{Solución de TSP con 1000 ciudades}
\end{figure}

Aunque esta vez se hace mucho mas dificultosa la tarea de entender y seguir el recorrido que la red encontró para atravesar todas las ciudades, si se observa con detenimiento se puede decir que logro el objetivo en ambos casos, sin embargo el tiempo demorado para la resolución del problema con $1000$ ciudades fue $4$ veces el tiempo demorado anteriormente.

Esto demuestra el potencial de este tipo de redes para rápidamente lograr una solución aproximada pero decente a problemas clásicos y computacionalmente complejos como lo es este. Cabe destacar que el hecho de que dicha aproximación se pueda lograr sin un algoritmo demasiado costoso, permite el uso de esta red para la solución de cualquier problema polinomialmente reducible a TSP.

\section{Reducción de dimensionalidad y descubrimiento de clusters}

En esta sección se aplicaran las redes de kohonen a tareas de reducción de la dimensionalidad, y a partir de esta reducción se graficara la \textit{matriz U} de la red para detectar agrupaciones de neuronas en vecindades que quedan lejanas a otras vecindades (es decir, clusters).

\subsection{Forma del dataset}

El conjunto de datos a utilizar cuenta con $500$ muestras de una variable con $100$ dimensiones. Luego de calcular los datos de resumen del conjunto se obtiene que, cada columna de una muestra es un numero entre $-1$ y $1$ y tiene una variabilidad considerablemente alta.

\begin{table}[H]
\centering
\small
\begin{tabular}{|c|c|c|c|}
\hline
\textbf{Dato} & \textbf{Col 0} & \textbf{Col 1} & \textbf{Col 2} \\
\hline
count & 500 & 500 & 500 \\
mean & -0.097 & -0.17 & -0.28 \\
std & 0.68 & 0.74 & 0.70 \\
min & -1.00 & -1.00 & -1.00 \\
25\% & -0.77 & -0.78 & -0.92 \\
50\% & -0.10 & -0.52 & -0.57 \\
75\% & 0.49 & 0.86 & 0.33 \\
max & 1.00 & 1.00 & 1.00 \\
\hline
\end{tabular}
\caption{Estadísticas descriptivas de las tres primeras columnas del dataset}
\end{table}

\subsection{Cálculo de la matriz \textit{U}}

La matriz \textit{U} permite visualizar regiones de alta o baja similitud en el mapa de Kohonen. Se construye calculando la distancia promedio entre cada neurona y sus vecinas adyacentes, lo cual permite detectar clusters y sus bordes.

El resultado es una matriz \( \mathbf{U} \) del mismo tamaño que el mapa, donde valores altos indican zonas de transición entre clusters y valores bajos indican regiones internamente homogéneas.

\subsection{Visualizaciones de matriz \textit{U} con diferentes redes}

A continuación se demuestran los resultados de dichas matrices con diferentes tamaños de redes de manera incremental. En los gráficos las zonas oscuras denotaran un valor bajo del promedio de distancias, lo que quiere decir que las neuronas en esa zona se activan ante estímulos parecidos, las zonas claras denotan que las neuronas que se encuentran ahí distan mucho de sus vecinas, lo que podría significar el borde entre clusters.

\begin{figure}[H]
\centering
\begin{minipage}[b]{0.20\textwidth}
  \centering
  \includegraphics[width=\textwidth]{img/um_15n10e.png}
  \caption*{(a) 15x15 neuronas, 10 épocas}
\end{minipage}
\hfill
\begin{minipage}[b]{0.20\textwidth}
  \centering
  \includegraphics[width=\textwidth]{img/um_30n10e.png}
  \caption*{(b) 30x30 neuronas, 10 épocas}
\end{minipage}

\vspace{0.3cm}

\begin{minipage}[b]{0.20\textwidth}
  \centering
  \includegraphics[width=\textwidth]{img/um_50n10e.png}
  \caption*{(c) 50x50 neuronas, 10 épocas}
\end{minipage}
\hfill
\begin{minipage}[b]{0.20\textwidth}
  \centering
  \includegraphics[width=\textwidth]{img/um_100n10e.png}
  \caption*{(d) 100x100 neuronas, 10 épocas}
\end{minipage}

\caption{Comparación de las matrices U generadas con distintos tamaños de mapa de Kohonen.}
\end{figure}

En la primera figura se pueden apreciar cuatro zonas oscuras bastante bien definidas y separadas por bordes claros, esto da la impresión de que efectivamente hay presente algún tipo de agrupación, pero no se puede decir con certeza. Luego la segunda figura, aunque también muestra cuatro zonas oscuras, unas son mucho mas predominantes que otras y no demuestra tan claramente agrupaciones. Sin embargo, a medida que se aumenta el tamaño del mapa, se termina descubriendo que efectivamente, hay clusters en el conjunto de datos y son un total de 4.

\section{Bonus: análisis de tendencia a clustering y k-means}

En esta sección y a modo de bonus, se investiga por otros medios la presencia de esos clusters descubiertos en la sección anterior.

\subsection{Silhouette y método del codo}

El método del codo consiste en aplicar el algoritmo \textit{k-means} para distintos valores de $k$ (número de clusters) y observar la variación de la suma de errores cuadráticos (inercia). A medida que $k$ aumenta, la inercia disminuye, pero a partir de cierto punto esa reducción se vuelve marginal. El punto de 'codo' sugiere el número óptimo de clusters.

El coeficiente de \textit{silhouette} mide qué tan bien se encuentra cada punto dentro de su clúster comparado con otros clusters. Toma valores entre $-1$ y $1$, donde valores cercanos a $1$ indican que el punto está bien asignado.

En este caso se calcularon ambos indices para valores de $k$ incrementales desde $2$ hasta $14$ y descubriendo que el valor posiblemente mas optimo, es el mismo valor anteriormente encontrado con las redes de kohonen y es $k=4$ como se puede visualizar a continuación.

\begin{figure}[H]
\centering
\includegraphics[width=0.49\textwidth]{img/elbow_silouette.png}
\caption{Método de codo y Silhouette}
\end{figure}

\subsection{Confirmación de clusters detectados con k-means}

Finalmente se ejecuta el algoritmo de k-Means con valores al rededor de la cantidad optima de clusters encontrada, y se utiliza PCA (análisis de componentes principales) para reducir la dimensionalidad y obtener gráficos de dos dimensiones donde se podrían observar los agrupamientos.

\begin{figure}[H]
\centering
\includegraphics[width=0.40\textwidth]{img/kmeans4.png}
\caption{K-means con $k=4$}
\end{figure}

En el gráfico se pueden observar a simple vista los 4 centroides de los grupos y los cuatro tonos de color que se utilizaron dependiendo de las etiquetas puestas por k-means. Esto re-confirma que la cantidad de grupos en este conjunto de datos es efectivamente la encontrada con las redes de kohonen y es cuatro.

A continuación se observan dos gráficos adicionales con numero de clusters $3$ y $5$, ya que son valores cercanos al optimo visto en el análisis con el método de codo y silhouette. En los mismos es fácil notar que los grupos no están perfectamente asignados ya que o hay un grupo que abarca demasiado (en el caso de $k=3$) o hay dos grupos disputando un centroide (en el caso de $k=5$).

\begin{figure}[H]
\centering
\includegraphics[width=0.40\textwidth]{img/kmeans3.png}
\caption{K-means con $k=3$}
\end{figure}

\begin{figure}[H]
\centering
\includegraphics[width=0.40\textwidth]{img/kmeans5.png}
\caption{K-means con $k=5$}
\end{figure}

\section*{Conclusión}

A lo largo de este trabajo se demostró la versatilidad y potencia de las redes auto-organizadas de Kohonen (SOM) en diversas aplicaciones. Primero, se verificó su capacidad para preservar la topología de distribuciones uniformes en figuras geométricas, evidenciando cómo las neuronas logran adaptarse a la forma del conjunto de datos mediante un aprendizaje no supervisado. Luego, se abordó el Problema del Viajante de Comercio, donde una red SOM unidimensional organizada en anillo logró aproximar recorridos eficientes para grandes cantidades de ciudades, validando su utilidad en problemas complejos. Finalmente, se aplicó una SOM a un conjunto de datos de alta dimensionalidad, utilizando la matriz U para revelar estructuras de clustering, que posteriormente fueron confirmadas con técnicas clásicas como k-means y análisis de silhouette. En conjunto, estos resultados destacan a las redes de Kohonen como una herramienta valiosa para la representación, reducción y organización de datos complejos.

\newpage

\begin{thebibliography}{9}

\bibitem{hertz}
Hertz, J., Krogh, A., \& Palmer, R. G. (1991).
\textit{Introduction to the Theory of Neural Computation}.
Addison-Wesley.

\bibitem{kohonen}
Kohonen, T. (1982).
Self-organized formation of topologically correct feature maps.
\textit{Biological Cybernetics}, 43(1), 59–69.

\end{thebibliography}
